\scanpage{30}
% @ 1–4
\uline{Corollaire} 1 \unsure{Tout|Pour} \deleted{\unsure{tout}} sous-module $M'$ de $M$, il y a un plus petit sous-module \inserted{$M'^{\unsure{\vcenter{\hbox{\tikz[x=1ex,y=1ex]{\draw[line width=.11ex,line cap=round,line join=round] (.57,.90) .. controls (.45,.76) and (.24,.53) .. (.15,.35) .. controls (.35,.35) and (.56,.45) .. (.73,.41) .. controls (.64,.2) and (.45,.03) .. (.32,-.14);}}}|\natural|\sharp|\mathrm{s}}}$} stable de $M$ qui \unsure{contienne|contient} $M'$, et c'est aussi \unsure{la|le} sous $\check A$-algèbre \unsure{engendrée|engendré}. Si $M'$ est de t.f. \deleted{\unsure{sur}} $k$, $M'^{\vcenter{\hbox{\tikz[x=1ex,y=1ex]{\draw[line width=.11ex,line cap=round,line join=round] (.57,.90) .. controls (.45,.76) and (.24,.53) .. (.15,.35) .. controls (.35,.35) and (.56,.45) .. (.73,.41) .. controls (.64,.2) and (.45,.03) .. (.32,-.14);}}}}$ aussi, \deleted{\unsure{donc}} \deleted{\unsure{M}}% end of the crossed group

% @ 5–6
\uline{Corollaire 2} $M$ est limite inductive \unsure{filtr.|filtrante}\deleted{\unsure{$\to$|t}} de ses sous-modules de t. f. \uline{stables}, \deleted{\unsure{par}}

% @ 7–8
\deleted{\uline{Corollaire 3}} \deleted{\unsure{M}} \deleted{Appliquons} Considérons $A$ comme \unsure{bicomodule|bimodule} sur lui-même à g. et à droite, \deleted{\unsure{plus}} donc comme bi-module sur $\check A$,
% @ 9–12
\deleted{Alors pour \unsure{tt} $x$, les \unsure{s.}} considérons les sous-bi-modules $B\subset A$. \deleted{On voit que pour} Appliquant Cor 1 et 2 \deleted{les sous-\unsure{bi-mod.}} \deleted{\unsure{à} $A$ considéré comme $A$} à \inserted{la coalgèbre produit ;} $A\otimes A^{\mathrm{o}}$ \inserted{($A^{\mathrm{o}}$ l'opposée de $A$)} et au module $A$ pour celle-ci, ou raisonnant directement, on trouve

% @ 13–16
\uline{Cor 3} \deleted{\unsure{Tout} \unsure{bi} \unsure{linéi}} Tout sous-bi-module $B$ de $A$ est contenu dans un plus-petit sous-bi-module $B^{\vcenter{\hbox{\tikz[x=1ex,y=1ex]{\draw[line width=.11ex,line cap=round,line join=round] (.57,.90) .. controls (.45,.76) and (.24,.53) .. (.15,.35) .. controls (.35,.35) and (.56,.45) .. (.73,.41) .. controls (.64,.2) and (.45,.03) .. (.32,-.14);}}}}$ \deleted{stable \unsure{par}} \deleted{\inserted{\unsure{qui est}}} stable à g. et à droite (savoir le sous-$\check A$-bi-module engendré), et si $B$ est de t. f. sur $k$, $B^{\vcenter{\hbox{\tikz[x=1ex,y=1ex]{\draw[line width=.11ex,line cap=round,line join=round] (.57,.90) .. controls (.45,.76) and (.24,.53) .. (.15,.35) .. controls (.35,.35) and (.56,.45) .. (.73,.41) .. controls (.64,.2) and (.45,.03) .. (.32,-.14);}}}}$ aussi.
% @ 17–18
Donc $A=\unsure{\underline{\lim}|\underrightarrow{\lim}}\,A_i$, les $A_i$ étant stables à g. et à droite.

% @ 19–23
Or cette stabilité \inserted{pour un $B$ donné} signifie que $\unsure{\pi|\Pi}:A\unsure{-|\to}A\otimes A$ envoie $B$ dans $(B\otimes A)\cap(A\otimes B)$. Si \deleted{\unsure{B est}} $A/B$ est plat sur $k$ (\unsure{p.~ex.} $k$ un corps) cel signifie aussi que $\pi$ envoie $B$ dans $B\otimes B$, ou encore qu'il y a une structure de co-algèbre induite. On trouve donc

% @ 24–26
\uline{Cor. 4} Toute coalgèbre \inserted{$A$} sur un corps est \uline{lim} \inserted{\unsure{fil.}} de coalgèbres de t. f., donc $\check A$ est $\underset{\leftarrow}{\mathrm{lim}}$ \inserted{\unsure{fil. stricte}} d'algèbres finies sur $k$.
% archive sheet number: 29
